\documentclass[10pt,a4paper]{amsart}
\usepackage[margin=19mm]{geometry}
\usepackage{amsmath,amssymb,mathtools}
\usepackage[scr=rsfs,cal=euler]{mathalfa}
\usepackage{xfrac}
\usepackage[unicode,hidelinks,bookmarks=false]{hyperref}
\newcommand{\ie}{i.e.\ }
\newcommand{\cf}{cf.\ }
\newcommand{\resp}{respectively\ }
\newcommand{\Subsection}{\subsection}
\providecommand{\nobreakdash}{\nobreak\hbox{-}}
\setlength{\emergencystretch}{2em}
\multlinegap0pt
\usepackage[shortlabels]{enumitem}
\usepackage{xcolor}
\usepackage{amssymb}
\usepackage{mathtools}
\usepackage{bm}
\usepackage[normalem]{ulem}
\usepackage{cancel}
\usepackage{tcolorbox}
\newtheorem{lemma}{Lemma}
\newcommand{\R}{\mathbb{R}}
\newcommand{\proj}{\pi}
\newcommand{\projX}{\proj_{\Omega}}
\title{The gap between a variational problem and its occupation measure relaxation}
\author{Milan Korda, Rodolfo Rios-Zertuche}
\date{Source: arXiv:2205.14132v4 (CC BY 4.0)}
\begin{document}
\maketitle

\noindent\emph{Source and licence.} The gap between a variational problem and its occupation measure relaxation. Version arXiv:2205.14132v4 (CC BY 4.0), distributed by arXiv under the Creative Commons Attribution 4.0 International licence, http://creativecommons.org/licenses/by/4.0/, as recorded in the arXiv licence metadata for this exact version and shown on the abstract page https://arxiv.org/abs/2205.14132v4. Full licence notice: \texttt{licenses/arxiv-2205.14132-CC-BY-4.0.txt}.\par\smallskip

\noindent\textit{Editorial scope.} Counted unit: the article's lemma lem:weakderivative (source0.tex lines 1335--1345). The article's complete original proof of it is delivered with it (source0.tex lines 1346--1502, one \texttt{proof} environment of 157 lines). No mathematics is written, repaired or re-derived: the statement and the proof are the article's own lines, in the article's own notation, and no retained line is substituted or re-flowed.  Retained uncounted as the source-local context the proof needs: lines 281--283; lines 1127--1137; lines 1174--1179; lines 1182--1184; lines 1192--1194; lines 1216--1223; lines 1227--1233.  One declared adaptation: exactly 3 ancillary strings inside the retained spans are omitted (image-inclusion commands and, where the source repeats a label, the repeated label definitions) and no other line differs from the source; the figure environments, with their placement, captions and labels, are kept, and the package records the omitted strings in fidelity receipts bound to the affected bodies.  No retained line contains a citation, so the wrapper carries no bibliography and no bibliography entry.  The retained lines contain the article's own diagram code, which the wrapper typesets with the standard tikz packages; no image file is delivered and the package declares no asset dependency.  Editorial formatting only, in addition: an AMS \texttt{amsart} preamble that restates the article's own declarations and macros used by the retained lines, namely the environments \texttt{lemma} on separate counters, together with the standard packages those lines need.  The retained environments print in this document as Lemma 1 in order of appearance, whereas the article prints the same items with the numbering of its own full document.  The complete native source of the article as distributed in the arXiv e-print for this exact version is bundled unchanged as \texttt{sources/126239/source0.tex}.
\begin{equation}\label{eq:defweakdiff}
 \int_\Omega \varphi \,D\phi\,dx=-\int_\Omega \phi \, D\varphi\,dx
\end{equation}

\begin{lemma}\label{lem:main}
The functions $\varphi_r:\Omega\to \R$ defined by,
\[\varphi_r(x)=\inf_{\substack{y\in Y\\\rho(x,y)\leq r}}y,\qquad r\in {\mathcal R}=\rho(\Omega\times Y)\subset(-\infty,0], \; x\in \Omega,\]
are weakly differentiable.
These functions satisfy, for all $X\in C^\infty_c(\Omega\times Y;\R^{n+1})$,
\begin{equation}\label{eq:divXvarphi}
    \int_{{\mathcal R}}\int_{\{(x,y)\in\Omega\times Y:y\geq \varphi_r(x)\}}\hspace{-20mm}\operatorname{div}X\,dx\,dy\,d\nu(r)=\int_{{\mathcal R}}\int_\Omega\langle X(x,\varphi_r(x)),(D\varphi_r(x),-1)\rangle\,dx\,d\nu(r),
\end{equation}
where $\nu$ is 
{the} Lebesgue measure restricted to ${\mathcal R}$.
\end{lemma}

 \begin{figure}
 
 \centering
 \caption{The sets $\rho^{-1}(-\infty,r]$ and their exterior unit normal $\eta_r$, as in the situation of Lemma \ref{lem:main}.}
  \label{fig:2} 
 \end{figure}

 \begin{equation}\label{eq:normal}
  \int_{\rho^{-1}(-\infty,r]}\operatorname{div}X\,dx\,dy=\int_{\partial \rho^{-1}(-\infty,r]} \langle X,\eta_r\rangle d H^n.
 \end{equation}

 \begin{equation}\label{eq:defzeta}
 [\zeta_r(x)]_i=-\frac{[\eta_r(x,\varphi_r(x))]_i}{[\eta_r(x,\varphi_r(x))]_{n+1}},\quad 1\leq i\leq n,
 \end{equation}

 \begin{lemma}\label{lem:function}
 For $r\in {\mathcal R}$, let $V_r$ be the set of points $(x,y)\in\Omega\times Y$ such that $\eta_r(x,y)$ is defined and its last coordinate is $[\eta_r(x,y)]_{n+1}=0$. For almost every $r\in {\mathcal R}$, we have 
 \[H^n(V_r)=0.\]
%
%
%
%
 \end{lemma}

  \begin{figure}
 
 \centering
 \caption{When there is a vertical segment $\{x\}\times[a,b]$ 
 in the boundary $\partial \rho^{-1}(-\infty,r]$, the normal vector is horizontal, that is, of the form $(z,0)$, $z\in\R^n$. The proof of Lemma \ref{lem:function} shows that the $n$-dimensional volume of the union of these segments is zero.}
  \label{fig:3}
 \end{figure}

 \begin{lemma}\label{lem:weakderivative}
 \begin{enumerate}[label=\roman*.,ref=(\roman*)]
 The following statements are true for almost every $r\in {\mathcal R}$:
  \item \label{it:weakderivative1}The vector field $\zeta_r$ defined in \eqref{eq:defzeta} is minus the weak derivative of $\varphi_r$, that is, 
  \[\int_\Omega\varphi_r(x)\nabla\phi(x)\,dx=\int_\Omega\phi(x)\zeta_r(x)\,dx\]
  for all $\phi\in C^\infty_c(\Omega)$.
  \item \label{it:weakderivative2} We also have, for $X\in C^1_c(\Omega\times Y;\R^{n+1})$, %
  \[\int_{\partial\rho^{-1}(-\infty,r]}\langle X,\eta_r\rangle dH^n
      =\int_{\Omega}\langle X(x,\varphi_r(x)),(\zeta_r(x),-1)\rangle dx.\]
  \end{enumerate}
 \end{lemma}

 \begin{proof}
 For almost every $r\in {\mathcal R}$, $\eta_r$ is well defined on almost all the boundary $\partial\rho^{-1}(-\infty,r]$, and by Lemma \ref{lem:function}, the set $V_r$ of points $(x,y)\in\partial\rho^{-1}(-\infty,r]$ with $[\eta_r(x,y)]_{n+1}=0$ has $n$-dimensional Hausdorff measure zero, $H^n(V_r)=0$. Let $r\in {\mathcal R}$ be a number that enjoys those properties.
 
 Denote by $m_1^r$ the Hausdorff measure $H^n$ on the boundary $(\Omega\times Y)\cap\partial \rho^{-1}(-\infty,r]$. The set $(\Omega\times Y)\cap\partial \rho^{-1}(-\infty,r]$ is depicted in Figure \ref{fig:2} this set is depicted as a thickened curve; since $\Omega$ is open, it does not include any points in $(\partial\Omega)\times Y$.
 %
 %
 %
 %
 Observe that $m_1^r$ also \emph{is equal to its restriction to the graph of $\varphi_r$,}
 \begin{equation}\label{eq:restriction}
  m^r_1=m^r_1|_{(\Omega\times Y)\cap\partial \rho^{-1}(-\infty,r]\setminus V_r}=m^r_1|_{\operatorname{graph}(\varphi_r)}
 \end{equation}
 because $H^n(V_r)=0$, and $V_r$ contains the part of $\partial\rho^{-1}(-\infty,r]$ corresponding to discontinuities of $\varphi_r$; see Figure \ref{fig:3}. 
 
 Denote by $m_2^r$ the pushforward of 
{the} Lebesgue measure on $\Omega$ by the map $\Psi(x)=(x,\varphi_r(x))$, that is, $m_2^r$ is a measure on $\Omega \times Y$ given by
 \[m_2^r=\Psi_{\#}dx|_{\Omega}.\] 
 
 The measure $m_2^r$ is absolutely continuous with respect to $m_1^r$. Indeed, if $A\subset\Omega\times Y$ is a measurable set of zero $m_1^r$ measure, this means that for every $\varepsilon>0$, $A$ can be covered with countably many balls $U_1,U_2,\dots$ such that $\sum_{i=1}^\infty(\operatorname{diam} U_i)^n\leq \varepsilon$, according to the definition of $n$-dimensional Hausdorff measure. The image $\projX{}(A)$ through the projection $\projX{}\colon\Omega\times Y\to\Omega$ can then be covered by the projections of the balls $\projX{}(U_i)$, which will still satisfy (for some $C>0$ dependent only on $n$)
 \[m^r_2(A)\leq m^r_2({\textstyle\bigcup_{i=1}^\infty U_i})\leq |\textstyle\bigcup_{i=1}^\infty \projX{}(U_i)|\leq C\sum_{i=1}^\infty(\operatorname{diam}\projX{}(U_i))^n=C\sum_{i=1}^\infty(\operatorname{diam}U_i)^n\leq \varepsilon,\] 
 and hence $A$ is a set of measure zero with respect to $m_2^r$. This proves the absolute continuity of $m_2^r$ with respect to $m_1^r$.
 
 By the Radon-Nikodym theorem, there is a measurable function $G_r(x,y)$ such that $m^r_2=G_rm^r_1$.
 The measure $m_2^r$ is also supported in $(\partial \rho^{-1}(-\infty,r])\cap(\Omega\times Y)$. For any measurable set $B\subset \Omega$, we have 
 \[\int_{\Omega\times Y}\chi_{B}(x)\,G_r(x,y)\,dm_1^r(x,y) =\int_{B\times Y}G_r(x,y)\,dm_1^r(x,y)=m_2^r(B\times Y)=|B|.\]
  This means that $(\projX)_\#(G_rm^r_1)$ is 
{the} Lebesgue measure on $\Omega$. From \eqref{eq:restriction} it follows that also $(\projX)_\#(G_rm^r_1|_{\operatorname{graph}(\varphi_r)})$ is 
{the} Lebesgue measure on $\Omega$, and this means that, for almost every $x\in \Omega$, $G_r(x,\varphi_r(x))>0$. Set
  \[J_r(x,y)=\begin{cases}1/G_r(x,y),&G_r(x,y)>0,\\0,&\text{else,}\end{cases}\]
  for $(x,y)\in \partial\rho^{-1}(-\infty,r]$; we then have $m_1^r=J_rm_2^r$.
 
 With $J_r$ defined like this, we have %
 for all measurable functions  $f\colon\Omega\times Y\to\R$, 
 \begin{multline}\label{eq:jacobian}
  \int_{\partial\rho^{-1}(-\infty,r]}f(x,y)\eta_r(x,y)\,dH^n=\int_{\partial\rho^{-1}(-\infty,r]}f(x,y)\eta_r(x,y)\,dm^r_1\\
  =\int_{\partial\rho^{-1}(-\infty,r]}f(x,y)\eta_r(x,y)J_r(x,y)\,dm_2^r\\=
  \int_\Omega f(x,\varphi_r(x))\eta_r(x,\varphi_r(x))J_r(x,\varphi_r(x))dx.
 \end{multline}

 From \eqref{eq:normal} it follows that
 \[\int_{\{y\geq \varphi_r(x)\}}\operatorname{div}X\,dx\,dy=\int_{\partial \rho^{-1}(-\infty,r]} \langle X,\eta_r\rangle d H^n,\quad X\in C^\infty_c(\overline{\Omega\times Y};\R^{n+1}),\]
 or equivalently, we have the vector-valued identity (proved from the above by taking $X=\phi \,e_i$ for each vector $e_i$ in the standard basis)
 \begin{equation}\label{eq:previouslydisplayedidentity}
 \int_{\{y\geq \varphi_r(x)\}}\nabla \phi(x,y)\,dx\,dy=\int_{\partial \rho^{-1}(-\infty,r]} \phi\, \eta_r\, d H^n,\quad \phi\in C^\infty_c(\overline{\Omega\times Y};\R).
 \end{equation}
%
%
%
%
%
%
%
%
%
%
%
%
%
 Take $N>0$ large enough that $\operatorname{supp}\mu\subset\overline\Omega\times[-N,N]$, and take a function $\psi\in C^\infty_c(\overline{\Omega\times Y})$, $0\leq\psi\leq 1$, such that $\psi(x,y)=1$ for all $|y|\leq N$. 
  Then, if we let $\mathbf n$ denote the unit normal to the boundary of $\Omega\times [0,+\infty)\subset\Omega\times Y$,  using \eqref{eq:jacobian}, and computing the derivatives below as in $\Omega\times Y$, so that for example $\nabla \phi(x)=\nabla_{x,y} \phi(x)=(\frac{\partial\phi}{\partial x_1},\dots,\frac{\partial\phi}{\partial x_n},0)$ to account for $\frac{\partial\phi}{\partial y}=0$, we have
 \begin{align*}
     \int_\Omega&\varphi_r(x)\nabla\phi(x)\,dx= \int_\Omega\int_0^{\varphi_r(x)}\nabla \phi(x)\,dy\,dx\\
     &=\int_\Omega\int_0^{\varphi_r(x)}\nabla (\psi\phi)(x,y)\,dy\,dx\\
     \notag &=\int_{\{(x,y)\in\Omega\times Y:0\leq y\leq \varphi_r(x)\}}\nabla(\psi\phi)\,dx\,dy
     -\int_{\{(x,y)\in\Omega\times Y:0\geq y\geq \varphi_r(x)\}}\nabla(\psi\phi)\,dx\,dy \\
     \notag &=\int_{\{(x,y)\in\Omega\times Y:y\geq 0\}}\nabla(\psi\phi)\,dx\,dy-\int_{\{(x,y)\in\Omega\times Y:y\geq \varphi_r(x)\}}\nabla(\psi\phi)\,dx\,dy \\
     \notag &=\int_{\partial\{(x,y)\in\Omega\times Y:y\geq 0\}} \psi\phi\mathbf{n}\,dH^n-\int_{\partial\rho^{-1}(-\infty,r]} \psi\phi\eta_r\,dH^n\\
     \notag &=\int_{\Omega\times \{0\}} \phi\mathbf{n}\,dH^n-\int_{\{(x,\varphi_r(x)):x\in\Omega\}} \phi\eta_r\,dH^n\\
     %
     &=\int_\Omega \phi(x)e_{n+1}dx+\int_\Omega\phi(x)\eta_r(x)J_r(x,\varphi_r(x))%
     dx
 \end{align*}
  for all $\phi\in C^\infty_c(\Omega;\R)$; here, we have first used the fact that $\varphi_r(x)=\int_0^{\varphi_r(x)}dy$. Then we introduced $\psi$, and we separated the positive and negative parts of $\varphi_r$. We have expressed them in a slightly different form (see Figure \ref{fig:4}), and then passed to the boundary using \eqref{eq:previouslydisplayedidentity} and its equivalent for the domain $\{y\geq 0\}\subset \overline{\Omega\times Y}$. In the next-to-last line we used that $\psi=1$ in the  domain of integration and we have kept only the parts that do not cancel out from the fact that $\phi$ is compactly supported in $\Omega$, namely, the sets $\Omega\times \{0\}\subset\partial\{y\geq 0\}$, whose normal is $e_{n+1}$,  and  $\varphi_r(\Omega)$, whose normal is $\eta_r$;  $J_r(x)$ comes in %
  once we apply \eqref{eq:jacobian}.
  %
  In other words, we have
  \begin{equation}\label{eq:vectident}
   \int_\Omega\varphi_r(x)\begin{pmatrix}\nabla_x \phi(x)\\0\end{pmatrix}dx
   =\int_\Omega\phi(x)e_{n+1}dx+\int_\Omega\phi(x)\eta_r(x)J_r(x,\varphi_r(x))dx,%
  \end{equation}
  where the 0 entry in the left-hand side appears because $\phi$ is independent of $y$.
  The last entry of \eqref{eq:vectident} gives
  %
  \[0=\int_\Omega\phi(x)1\,dx+\int_\Omega \phi(x)[\eta_r(x,\varphi_r(x))]_{n+1}J_r(x,\varphi_r(x))\,dx.\]
  %
  Since this is true for all $\phi\in C^\infty_c(\Omega)$, we conclude that, for almost every $x\in\Omega$, \[J_r(x,\varphi_r(x))=-\frac1{[\eta_r(x,\varphi_r(x))]_{n+1}},\]
  and \eqref{eq:jacobian} becomes (cf.~\eqref{eq:defzeta})
  \begin{equation}\label{eq:jacobian2}
  \int_{\partial\rho^{-1}(-\infty,r]}f(x,y)\eta_r(x,y)\,dH^n(x,y)=\int_{\Omega}f(x,\varphi_r(x))\
  \begin{pmatrix}
   \zeta_r(x)\\-1
  \end{pmatrix}dx.
  \end{equation}
  Applying \eqref{eq:jacobian2} to $f=X_i$ to obtain, in the $i$-th entry,
  \begin{equation*}%
  \int_{\partial\rho^{-1}(-\infty,r]}X_i(x,y)[\eta_r(x,y)]_i\,dH^n(x,y)=\begin{cases}\int_{\Omega}X_i(x,\varphi_r(x))
   [\zeta_r(x)]_i\,dx,& 1\leq i\leq n-1,\\
   -\int_{\Omega}X_n(x,\varphi_r(x))
   \,dx,&i=n,
   \end{cases}
  \end{equation*}  
  and adding these over all $i=1,\dots, n$ proves the identity in item \ref{it:weakderivative2}.
 
%
%
%
%
%
%
%
%
%
%
%
%
%
%
%
%
%
%
%
%
%
%
%
%
%
%
%
%
%
%
%
%
%
%
%
%
%
%
%
%
  
    \begin{figure}
 
 \centering
 \caption{Illustrating a step in the proof of Lemma \ref{lem:weakderivative}, we see that the difference of integrals on the shaded areas in the first two diagrams on the left is equal to the difference of integrals on the two  areas on the right. %
 }
  \label{fig:4}
 \end{figure}
  
  We also have, taking only the first $n$ entries in \eqref{eq:vectident},
  \[\int_\Omega\varphi_r(x)\nabla\phi(x)\,dx=\int_\Omega\phi(x)\zeta_r(x)\,dx.\]
  This is precisely the definition \eqref{eq:defweakdiff} of weak differentiability.
%
\end{proof}

\end{document}
